% 底本：EGAIV.1.pdf，PDF第166--168页，原分册印刷第167--169页。
\subsection{扩张的可容许域}
\egap{21.6.1}
给定四个域 $k_0\subseteq k\subseteq K\subseteq L$，
由（20.6.16）和（20.6.17）有正合列
\begin{equation}
 0\longrightarrow\Upsilon_{K/k/k_0}\otimes_KL
 \xrightarrow{v}\Upsilon_{L/k/k_0}
 \xrightarrow{u}\Upsilon_{L/K/k_0}
 \xrightarrow{s}\Upsilon_{L/K/k}\longrightarrow0.
 \tag{21.6.1.1}
\end{equation}
固定 $k_0,K,L$，让中间域 $k$ 在 $k_0$ 和 $K$ 之间变化。
当 $k=k_0$ 时，当然有 $\Upsilon_{L/K/k_0}=\Upsilon_{L/K/k}$。
一般地，若（21.6.1.1）中的典范同态 $s$ 仍为双射，
就称 $k$ 是扩张 $L/K$ 的一个\textbf{$k_0$-可容许域}。
在一定条件下能找到这样的 $k$，而且它又“充分接近”$K$（例如 $[K:k]$ 有限），
这一点很有用：在某些问题中，它使我们能够把可能“过大”的微分模
$\Omega^1_{K/k_0}$ 和 $\Omega^1_{L/k_0}$（例如 $k_0$ 为素域时），
换成较易处理的 $\Omega^1_{K/k}$ 和 $\Omega^1_{L/k}$。
若 $k_0$ 为素域，简称“可容许域”。

置
\begin{equation}
\begin{split}
 \Delta(L/K,k/k_0)
 & =\Coker\bigl(\Upsilon_{K/k/k_0}\otimes_KL
                    \longrightarrow\Upsilon_{L/k/k_0}\bigr)\\
 & \simeq\Ker\bigl(\Upsilon_{L/K/k_0}
                    \longrightarrow\Upsilon_{L/K/k}\bigr).
\end{split}
\tag{21.6.1.2}
\end{equation}
这是一个 $L$-向量空间。它的维数记为 $d(L/K,k/k_0)$，
称为 $k$ 对扩张 $L/K$ 的\textbf{$k_0$-可容许性缺陷}；
它为零恰好意味着 $k$ 对这一扩张是 $k_0$-可容许的。
若 $k_0$ 为素域，则省略它，写作 $\Delta(L/K,k)$ 和 $d(L/K,k)$。

\egas{21.6.2}{命题}
设 $k_0\subseteq k\subseteq K\subseteq L$ 为四个域。

（i）下列条件等价：
\begin{enumerate}[label=\alph*),leftmargin=2.5em]
\item $k$ 对扩张 $L/K$ 是 $k_0$-可容许的，即
$\Upsilon_{L/K/k_0}\to\Upsilon_{L/K/k}$ 单射，因而双射；
\item 典范同态 $u:\Upsilon_{L/k/k_0}\to\Upsilon_{L/K/k_0}$ 为零；
\item 典范同态 $v:\Upsilon_{K/k/k_0}\otimes_KL\to\Upsilon_{L/k/k_0}$ 满射，因而双射；
\item $d(L/K,k/k_0)=0$，即 $\Delta(L/K,k/k_0)=0$。
\end{enumerate}
（ii）在下列任一情形中，（i）的等价条件成立：
$\alpha$）$L/k$ 可分；$\beta$）$L/K$ 可分；
$\gamma$）$k\subseteq k_0(K^p)$，其中 $p$ 表示 $k_0$ 的特征指数。

\eproof
（i）直接来自（21.6.1.1）的正合性。
（ii）若 $L/k$ 可分，则 $\Upsilon_{L/k/k_0}=0$（20.6.19），故（i）b)成立。
若 $L/K$ 可分，则 $\Upsilon_{L/K/k_0}=0$（20.6.19），故（i）a)成立。
最后，若 $k\subseteq k_0(K^p)$，则由（21.1.5.1）
$\Upsilon_{L/K/k}=\Upsilon_{L/K/k_0}$，故（i）a)成立。
\cmt{L12-N02}{这里及下文“特征指数”在特征零时取 $p=1$，在正特征时取相应素数。
因此特征零的 $\gamma$）自动满足；这与所有域扩张在特征零时均可分一致。
不要把“指数 $p=1$”误写成“特征为1”。}
\finishproof

\egap{21.6.3}
设有域单同态组成的交换图
\[
\begin{tikzcd}[column sep=large]
 k'_0\arrow[r]&k'\arrow[r]&K'\arrow[r]&L'\\
 k_0\arrow[u]\arrow[r]&k\arrow[u]\arrow[r]&K\arrow[u]\arrow[r]&L\arrow[u]
\end{tikzcd}
\]
由（20.6.21，ii）的自然性，得到典范同态
\begin{equation}
 \Delta(L/K,k/k_0)\longrightarrow\Delta(L'/K',k'/k'_0).
 \tag{21.6.3.1}
\end{equation}
\cmt{L12-T03}{原文引（20.6.17，ii），但该条没有（ii）这一分项。
这里使用的是（20.6.21，ii）中正合列的自然性及（20.6.4）中的函子性。
两边的标量域分别为 $L,L'$，该典范映射应按 $L\to L'$ 理解为双同态，不能无条件称为同一标量域上的线性映射。}
这些同态具有复合相容性，故 $\Delta(L/K,k/k_0)$ 可看作四元组 $(k_0,k,K,L)$ 的函子。

\egas{21.6.4}{命题}
（i）设 $k\subseteq k'\subseteq k''\subseteq K\subseteq L$ 为五个域。
则有由典范同态组成的正合列
\begin{equation}
\begin{split}
 0\longrightarrow\Delta(L/K,k'/k)
 \longrightarrow\Delta(L/K,k''/k)\\
 \longrightarrow\Delta(L/K,k''/k')\longrightarrow0,
\end{split}\tag{21.6.4.1}
\end{equation}
因而
\begin{equation}
 d(L/K,k''/k)=d(L/K,k''/k')+d(L/K,k'/k).
 \tag{21.6.4.2}
\end{equation}
（ii）设 $k_0\subseteq k\subseteq K\subseteq L\subseteq M$ 为五个域。
则有由典范同态组成的正合列
\begin{equation}
\begin{split}
 0\longrightarrow\Delta(L/K,k/k_0)\otimes_LM
 \longrightarrow\Delta(M/K,k/k_0)\\
 \longrightarrow\Delta(M/L,k/k_0)\longrightarrow0,
\end{split}\tag{21.6.4.3}
\end{equation}
因而
\begin{equation}
 d(M/K,k/k_0)=d(M/L,k/k_0)+d(L/K,k/k_0).
 \tag{21.6.4.4}
\end{equation}

\eproof
（i）考虑交换图
\[
\begin{tikzcd}[column sep=small,row sep=large,cramped]
 0\arrow[r]&\Upsilon_{K/k'/k}\otimes_KL\arrow[r]\arrow[d]
 &\Upsilon_{K/k''/k}\otimes_KL\arrow[r]\arrow[d]
 &\Upsilon_{K/k''/k'}\otimes_KL\arrow[r]\arrow[d]&0\\
 0\arrow[r]&\Upsilon_{L/k'/k}\arrow[r]\arrow[d]
 &\Upsilon_{L/k''/k}\arrow[r]\arrow[d]
 &\Upsilon_{L/k''/k'}\arrow[r]\arrow[d]&0\\
 0\arrow[r]&\Delta(L/K,k'/k)\arrow[r]
 &\Delta(L/K,k''/k)\arrow[r]
 &\Delta(L/K,k''/k')\arrow[r]&0.
\end{tikzcd}
\]
将三行分别看作复形 $T_1^\bullet,T_2^\bullet,T_3^\bullet$。
正合列（21.6.1.1）及定义（21.6.1.2）给出逐次数的正合列
$0\to T_1^\bullet\to T_2^\bullet\to T_3^\bullet\to0$。
由（21.6.1.1），$T_1^\bullet$ 与 $T_2^\bullet$ 的上同调除同一个次数外均为零，
而在该次数两者都典范地识别为
$\Upsilon_{k''/k'/k}\otimes_{k''}L$；
且 $T_1^\bullet\to T_2^\bullet$ 在此识别下诱导恒等映射。
上同调长正合列遂表明 $T_3^\bullet$ 的上同调全为零，得到（i）。
\cmt{L12-G04}{原文以“两复形有相同上同调”直接推出第三个无上同调。
所需的不只是两个上同调对象抽象同构，而是所给复形映射诱导的同态确为同构。
这里两者的核都经（21.6.1.1）识别为 $\Upsilon_{k''/k'/k}\otimes_{k''}L$，
自然性使诱导映射为恒等，正文补出这一点；命题结论不变。}

（ii）同样考虑交换图
\[
\begin{tikzcd}[column sep=small,row sep=large,cramped]
 0\arrow[r]&\Delta(L/K,k/k_0)\otimes_LM\arrow[r]\arrow[d]
 &\Delta(M/K,k/k_0)\arrow[r]\arrow[d]
 &\Delta(M/L,k/k_0)\arrow[r]\arrow[d]&0\\
 0\arrow[r]&\Upsilon_{L/K/k_0}\otimes_LM\arrow[r]\arrow[d]
 &\Upsilon_{M/K/k_0}\arrow[r]\arrow[d]
 &\Upsilon_{M/L/k_0}\arrow[r]\arrow[d]&0\\
 0\arrow[r]&\Upsilon_{L/K/k}\otimes_LM\arrow[r]
 &\Upsilon_{M/K/k}\arrow[r]
 &\Upsilon_{M/L/k}\arrow[r]&0.
\end{tikzcd}
\]
将三行看作 $T_1'^\bullet,T_2'^\bullet,T_3'^\bullet$，
（21.6.1.1）和（21.6.1.2）给出正合列
$0\to T_1'^\bullet\to T_2'^\bullet\to T_3'^\bullet\to0$。
由（21.6.1.1），$T_2'^\bullet$、$T_3'^\bullet$ 只有同一个次数的上同调可能非零，
两者都典范地识别为 $\Upsilon_{M/L/K}$，且诱导映射为恒等。
因此上同调长正合列给出 $T_1'^\bullet$ 无上同调，证明（ii）。
\finishproof

\egas{21.6.5}{推论}
（i）给定五个域 $k\subseteq k'\subseteq k''\subseteq K\subseteq L$，
$k''$ 对扩张 $L/K$ 是 $k$-可容许的，当且仅当 $k'$ 对该扩张是 $k$-可容许的，
且 $k''$ 对该扩张是 $k'$-可容许的。

（ii）给定五个域 $k_0\subseteq k\subseteq K\subseteq L\subseteq M$，
$k$ 对扩张 $M/K$ 是 $k_0$-可容许的，当且仅当它对扩张 $L/K$ 及 $M/L$
都为 $k_0$-可容许的。

这些断言由（21.6.4.2）和（21.6.4.4）立即得到，因为 $d$ 的值非负。

\egas{21.6.6}{推论}
设 $k_0\subseteq k\subseteq K\subseteq L$ 为四个域，
且 $k$ 对扩张 $L/K$ 是 $k_0$-可容许的。
若另有域 $k'_0,k',K',L'$，满足
\[
 k_0\subseteq k'_0\subseteq k'\subseteq k\subseteq K
 \subseteq K'\subseteq L'\subseteq L,
\]
则 $k'$ 对扩张 $L'/K'$ 是 $k'_0$-可容许的。
\cmt{L12-G05}{此处原文没有另列证明。
可先用（21.6.4，ii）的缺陷加法性，将 $L/K$ 缩至 $L'/K'$，其非负缺陷仍为零；
再用（21.6.4，i）沿 $k_0\subseteq k'_0\subseteq k'\subseteq k$ 分解基域缺陷，
每项都非负，故 $d(L'/K',k'/k'_0)=0$。这里只补充该推论与前条之间的推理联系。}
